% Options for packages loaded elsewhere
% Options for packages loaded elsewhere
\PassOptionsToPackage{unicode}{hyperref}
\PassOptionsToPackage{hyphens}{url}
\PassOptionsToPackage{dvipsnames,svgnames,x11names}{xcolor}
%
\documentclass[
  ngerman,
]{agujournal2019}
\usepackage{xcolor}
\usepackage{amsmath,amssymb}
\setcounter{secnumdepth}{5}
\usepackage{iftex}
\ifPDFTeX
  \usepackage[T1]{fontenc}
  \usepackage[utf8]{inputenc}
  \usepackage{textcomp} % provide euro and other symbols
\else % if luatex or xetex
  \usepackage{unicode-math} % this also loads fontspec
  \defaultfontfeatures{Scale=MatchLowercase}
  \defaultfontfeatures[\rmfamily]{Ligatures=TeX,Scale=1}
\fi
\usepackage{lmodern}
\ifPDFTeX\else
  % xetex/luatex font selection
\fi
% Use upquote if available, for straight quotes in verbatim environments
\IfFileExists{upquote.sty}{\usepackage{upquote}}{}
\IfFileExists{microtype.sty}{% use microtype if available
  \usepackage[]{microtype}
  \UseMicrotypeSet[protrusion]{basicmath} % disable protrusion for tt fonts
}{}
\makeatletter
\@ifundefined{KOMAClassName}{% if non-KOMA class
  \IfFileExists{parskip.sty}{%
    \usepackage{parskip}
  }{% else
    \setlength{\parindent}{0pt}
    \setlength{\parskip}{6pt plus 2pt minus 1pt}}
}{% if KOMA class
  \KOMAoptions{parskip=half}}
\makeatother
% Make \paragraph and \subparagraph free-standing
\makeatletter
\ifx\paragraph\undefined\else
  \let\oldparagraph\paragraph
  \renewcommand{\paragraph}{
    \@ifstar
      \xxxParagraphStar
      \xxxParagraphNoStar
  }
  \newcommand{\xxxParagraphStar}[1]{\oldparagraph*{#1}\mbox{}}
  \newcommand{\xxxParagraphNoStar}[1]{\oldparagraph{#1}\mbox{}}
\fi
\ifx\subparagraph\undefined\else
  \let\oldsubparagraph\subparagraph
  \renewcommand{\subparagraph}{
    \@ifstar
      \xxxSubParagraphStar
      \xxxSubParagraphNoStar
  }
  \newcommand{\xxxSubParagraphStar}[1]{\oldsubparagraph*{#1}\mbox{}}
  \newcommand{\xxxSubParagraphNoStar}[1]{\oldsubparagraph{#1}\mbox{}}
\fi
\makeatother


\usepackage{longtable,booktabs,array}
\newcounter{none} % for unnumbered tables
\usepackage{calc} % for calculating minipage widths
% Correct order of tables after \paragraph or \subparagraph
\usepackage{etoolbox}
\makeatletter
\patchcmd\longtable{\par}{\if@noskipsec\mbox{}\fi\par}{}{}
\makeatother
% Allow footnotes in longtable head/foot
\IfFileExists{footnotehyper.sty}{\usepackage{footnotehyper}}{\usepackage{footnote}}
\makesavenoteenv{longtable}
\usepackage{graphicx}
\makeatletter
\newsavebox\pandoc@box
\newcommand*\pandocbounded[1]{% scales image to fit in text height/width
  \sbox\pandoc@box{#1}%
  \Gscale@div\@tempa{\textheight}{\dimexpr\ht\pandoc@box+\dp\pandoc@box\relax}%
  \Gscale@div\@tempb{\linewidth}{\wd\pandoc@box}%
  \ifdim\@tempb\p@<\@tempa\p@\let\@tempa\@tempb\fi% select the smaller of both
  \ifdim\@tempa\p@<\p@\scalebox{\@tempa}{\usebox\pandoc@box}%
  \else\usebox{\pandoc@box}%
  \fi%
}
% Set default figure placement to htbp
\def\fps@figure{htbp}
\makeatother


% definitions for citeproc citations
\NewDocumentCommand\citeproctext{}{}
\NewDocumentCommand\citeproc{mm}{%
  \begingroup\def\citeproctext{#2}\cite{#1}\endgroup}
\makeatletter
 % allow citations to break across lines
 \let\@cite@ofmt\@firstofone
 % avoid brackets around text for \cite:
 \def\@biblabel#1{}
 \def\@cite#1#2{{#1\if@tempswa , #2\fi}}
\makeatother
\newlength{\cslhangindent}
\setlength{\cslhangindent}{1.5em}
\newlength{\csllabelwidth}
\setlength{\csllabelwidth}{3em}
\newenvironment{CSLReferences}[2] % #1 hanging-indent, #2 entry-spacing
 {\begin{list}{}{%
  \setlength{\itemindent}{0pt}
  \setlength{\leftmargin}{0pt}
  \setlength{\parsep}{0pt}
  % turn on hanging indent if param 1 is 1
  \ifodd #1
   \setlength{\leftmargin}{\cslhangindent}
   \setlength{\itemindent}{-1\cslhangindent}
  \fi
  % set entry spacing
  \setlength{\itemsep}{#2\baselineskip}}}
 {\end{list}}
\usepackage{calc}
\newcommand{\CSLBlock}[1]{\hfill\break\parbox[t]{\linewidth}{\strut\ignorespaces#1\strut}}
\newcommand{\CSLLeftMargin}[1]{\parbox[t]{\csllabelwidth}{\strut#1\strut}}
\newcommand{\CSLRightInline}[1]{\parbox[t]{\linewidth - \csllabelwidth}{\strut#1\strut}}
\newcommand{\CSLIndent}[1]{\hspace{\cslhangindent}#1}

\ifLuaTeX
\usepackage[bidi=basic,shorthands=off]{babel}
\else
\usepackage[bidi=default,shorthands=off]{babel}
\fi
\ifLuaTeX
  \usepackage{selnolig} % disable illegal ligatures
\fi


\setlength{\emergencystretch}{3em} % prevent overfull lines

\providecommand{\tightlist}{%
  \setlength{\itemsep}{0pt}\setlength{\parskip}{0pt}}



 


\usepackage{url} %this package should fix any errors with URLs in refs.
\usepackage{lineno}
\usepackage[inline]{trackchanges} %for better track changes. finalnew option will compile document with changes incorporated.
\usepackage{soul}
\linenumbers
\makeatletter
\@ifpackageloaded{caption}{}{\usepackage{caption}}
\AtBeginDocument{%
\ifdefined\contentsname
  \renewcommand*\contentsname{Inhaltsverzeichnis}
\else
  \newcommand\contentsname{Inhaltsverzeichnis}
\fi
\ifdefined\listfigurename
  \renewcommand*\listfigurename{Abbildungsverzeichnis}
\else
  \newcommand\listfigurename{Abbildungsverzeichnis}
\fi
\ifdefined\listtablename
  \renewcommand*\listtablename{Tabellenverzeichnis}
\else
  \newcommand\listtablename{Tabellenverzeichnis}
\fi
\ifdefined\figurename
  \renewcommand*\figurename{Abbildung}
\else
  \newcommand\figurename{Abbildung}
\fi
\ifdefined\tablename
  \renewcommand*\tablename{Tabelle}
\else
  \newcommand\tablename{Tabelle}
\fi
}
\@ifpackageloaded{float}{}{\usepackage{float}}
\floatstyle{ruled}
\@ifundefined{c@chapter}{\newfloat{codelisting}{h}{lop}}{\newfloat{codelisting}{h}{lop}[chapter]}
\floatname{codelisting}{Listing}
\newcommand*\listoflistings{\listof{codelisting}{Listingverzeichnis}}
\makeatother
\makeatletter
\makeatother
\makeatletter
\@ifpackageloaded{caption}{}{\usepackage{caption}}
\@ifpackageloaded{subcaption}{}{\usepackage{subcaption}}
\makeatother
\usepackage{bookmark}
\IfFileExists{xurl.sty}{\usepackage{xurl}}{} % add URL line breaks if available
\urlstyle{same}
% fallback for those not using the hyperref driver hyperxmp:
\makeatletter
\@ifundefined{xmpquote}{\newcommand{\xmpquote}[1]{#1}}{}
\makeatother
\hypersetup{
  pdftitle={Stadt als geopolitisches Werkzeug - das Beispiel Belgrad},
  pdfauthor={Henning Füller; Gruppe Studenti pobeđuju},
  pdflang={de},
  pdfkeywords={\xmpquote{Geopolitik}, \xmpquote{Everyday Geopolitics}},
  colorlinks=true,
  linkcolor={blue},
  filecolor={Maroon},
  citecolor={Blue},
  urlcolor={Blue},
  pdfcreator={LaTeX via pandoc}}


\journalname{Manuskript für Zeitschrift sub\textbackslash{}urban}

\draftfalse

\begin{document}
\title{Stadt als geopolitisches Werkzeug - das Beispiel Belgrad}

\authors{Henning Füller\affil{1}, Gruppe Studenti pobeđuju\affil{1}}
\affiliation{1}{Humboldt-Universität zu Berlin, }
\correspondingauthor{Henning Füller}{henning.fueller@geo.hu-berlin.de}


\begin{abstract}
Der Beitrag macht ausgehend von den beiden empirischen Beispielen
Waterfront und EXPO 2027 deutlich, inwiefern Versuche der geopolitischen
Positionierung regelmäßig auf räumliche Stabilisierungen angewiesen sind
und inwiefern diese in Teilen geopolitisch motivierte Fixierung
städtische Wirklichkeit prägt, aber auch neue Möglichkeiten für soziale
Bewegung schafft.
\end{abstract}

\section*{Plain Language Summary}
Wie wird Stadt zu einem geopolitischen Werkzeug und welche Möglichkeiten
eröffnet das für soziale Bewegungen?




\section{Einleitung}\label{einleitung}

Noch am Wahlabend des 25. Oktober 2026 hat Aleksandar Vučić,
Staatspräsident von Serbien seit 2014, seinen Rücktritt bekannt gegeben.
Die Studierendenliste kam auf 51,5\% der abgegebenen Stimmen. Vučićs SNS
nur auf 40,4 \%. Damit ist ein Staatschef zurückgetreten, den Beobachter
als ‚den von Europa meist favorisierten Autokraten` (Eror, 2018)
bezeichnet haben. Seit dem damals ebenfalls überraschenden Wahlsieg
seiner ‚Serbischen Fortschrittspartei` (SNS) im Jahr 2012 konnte Vučić
seine autoritäre Herrschaft immer weiter konsolidieren. Widerstand
begann sich vor allem seit November 2024 breit zu formieren. Damals war
ein Teils des Bahnhofs von Novi Sad eingestürzt und 16 Menschen sind auf
tragische Weise ums Leben gekommen. Das Ereignis löste Proteste und eine
zivilgesellschaftliche Gegenbewegung aus. Die Wahl und der Sieg der
Studierendenliste markiert nun einen unerwarteten Erfolg dieser
Bewegung.

Ausgehend von einem einjährigen Studienprojekt in Kooperation mit der
Universität Belgrad und einem einwöchigen Forschungsaufenthalt im Juni
2026 gehen wir in diesem Aufsatz einigen Hintergründen nach. Unserer
Lesart nach lässt sich die Regierung Vučić als paradigmatisches Beispiel
für einen „autoritären Neoliberalismus`` fassen (Bruff, 2014). Bereits
gut herausgearbeitet, auch in dieser Zeitschrift, ist die besondere
Rolle des Städtischen für dieses Regime. Städtische Großprojekte wie
z.B. die umfassende Neugestaltung des Save Ufers (‚Belgrade Waterfront`)
wurden zu einem Dreh- und Angelpunkt für die Legitmation und den
Machterhalt der Regierung Vučić (Čamprag, 2023). Solche Großprojekte
eröffnen unterschiedliche Möglichkeiten für Machterhalt. Sie versehen
das Regime mit einem Nimbus von Entwicklung und Modernisierung und sie
verschaffen Spielräume um lokale Eliten mit Investitionsmöglichkeiten
und Gewinnbeteiligung zu ko-optieren. Nicht zuletzt können durch solche
Großprojekte auch geopolitische Verbindungen gestärkt werden (Darieva
u.~a., 2026).

Im Detail betrachtet verweisen die städtischen Projekte also auf eine
neue geopolitische Skalenpolitik, die kennzeichnend ist für diese Form
des autoritären Neoliberalismus (Piletić, 2022). Einerseits kommt es zu
einer neuen Zentralisierung von Planungs- und Entscheidungskompetenzen
auf der nationalen Ebene. Andererseits bekommen sowohl lokale
wirtschaftliche Akteure aber auch internationale Investoren neue
Spielräume. Auf diesem Weg werden konkrete geopolitische Strategien
relevant für die lokale Stadtentwicklung.

Lokale Projekte sind dann über die Möglichkeit von Profit und
Klientelismus hinau auch ein Medium außenpolitischer Bündnispflege
(Golf-Kapital bei der Waterfront, die EXPO als internationale Bühne).
Aus dieser Doppelfunktion ergibt sich eine besondere Verwundbarkeit für
Protest, so die Hypothese des Beitrags.

\section{}\label{section}

Machterhalt wurde in diesem Fall entscheidend auf eine Skalenpolitik
gestützt, die insbesondere städtische Großprojekte zum Dreh- und
Angelpunkt genommen hat (Piletić, 2022). Städtische Politiken und
geopolitische Strategien kommen dabei in direkte Wechselwirkung.
Ergänzend zu diesen Befunden der Literatur schlagen wir im Folgenden
vor, die

Die ‚urban geopolitical turn` anhand zweier exemplarischer Fälle.
Empirisches Material sind neben teilnehmender Beobachtung und
Expert*inneninterviews mehr als 30 Stunden Interviewmaterial mit
Akteur*innen unterschiedlicher sozialer Bewegungen im Rahmen der
Studierendenproteste.

Die spezifische Form autoritärer Herrschaft auf dem Westbalkan wird
häufig als ‚Stabilokratie` diskutiert. Das Versprechen auf Stabilität
und rechtmäßige Wahlen wird nach außen stark gemacht und so das
Wohlwollen z.B. der EU sichergestellt. Gleichzeitig kommt es nach innen
zu starken autoritären Tendenzen und einer Aushöhlung von Demokratie und
Rechtsstaatlichkeit. Insbesondere das Regime von Aleksandar Vučić greift
zudem auf ökonomische Großprojekte mit intransparenter internationaler
Finanzierung als Mittel zur Stabilisierung nach innen zurück (Bieber,
2020). Die Stadt wird zu einem Werkzeug geopolitischer Positionierung.
Das Beispiel Belgrade Waterfront ist hier besonders augenfällig (Zeković
u.~a., 2023). Hochpreisige Apartmenttürme gebaut mit internationalem
Investmentkapital sollen den Anspruch Belgrads als \emph{world class
city} auch baulich zementieren. Dazu wurden allerdings bestehende
städtische Planungs- und Vergabeverfahren mit nationalem Recht
überschrieben (Grubbauer \& Čamprag, 2019). Die offensichtliche
Vereinnahmung zentraler Flächen der Stadt als Beute von
Investmentstrategien sowie die generische Architektur und mangelhafte
bauliche Qualität sind im Gegenzug Katalysator für die Formierung
zivilgesellschaftlichen Widerstands (Čamprag, 2024). Eine zweites
Fallbeispiel ist die Planung der EXPO 2027, die sich in ganz ähnlicher
Weise als geopolitisches Werkzeug und als Inkubator von Protest lesen
lässt.

Der Beitrag macht ausgehend von den beiden empirischen Beispielen
Waterfront und EXPO 2027 deutlich, inwiefern Versuche der geopolitischen
Positionierung regelmäßig auf räumliche Stabilisierungen angewiesen sind
und inwiefern diese in Teilen geopolitisch motivierte Fixierung
städtische Wirklichkeit prägt, aber auch neue Möglichkeiten für soziale
Bewegung schafft. ¡

\section{Stadt als Werkzeug im autoritären
Neoliberalismus}\label{sec-hintergrund}

Die Position auf der Grenze geopolitischer Bruchlinien ist ein
historisch prägendes Merkmal für Serbien. Solche Bruchlinien verliefen
zwischen westlichem Katholizismus und und orthodoxem Christentum,
zwischen Habsburger Monarchie und osmanischem Reich, zwischen
jugoslawischem Sozialismus und serbischem Nationalismus und heute
zwischen Europa und Russland. Die fortlaufend notwendige Aushandlungen
solcher geopolitischer Selbstpositionierung steht in einem besonderen
Spannungsverhältnis zur Stadtentwicklung, insbesondere in Belgrad. Die
‚ontologische Unsicherheit` bezüglich der geopolitischen Situation macht
die Stabilisierung mit Hilfe greifbarer baulicher Merkmale hier
besonders relevant Cupać (2020). Die Rolle des Städtischen im
Wechselspiel von staatlicher Macht, geopolitischen Diskursen und
Raumproduktion ist bisher zu vage bestimmt (Bădescu, 2022). Herrschende
Kräfte greifen aktiv auf gebauten Raum zu, auf funktionale und auf
symbolische Räume der Stadt, um bestimmte Positionierungen abzusichern
und auf Dauer zu stellen. Zugleich bieten solche städtische Strukturen
geeignete Ansatz- und Kristallisationspunkte für Widerstand.

\begin{quote}
a number of prominent contributions have drawn attention to the
important role of urban space in anchoring authoritarian
neoliberalism(s) (e.g.~Di Giovanni, 2016; Jenss, 2019; Tansel, 2019)
\end{quote}

Großprojekte der Stadtentwicklung sind für das System des ‚autoritären
Neoliberalismus` (Čamprag, 2023) in Serbien in vielfacher Hinsicht
funktional gewesen. Insbesondere kommt in Großprojekten eine
charakteristische Skalenpolitik zum Ausdruck, die Machtverhältnisse
zwischen internationaler, nationaler und städtischer Ebene neu ordnet
(Piletić, 2022). Waterfront

\section{Expo 2027}\label{expo-2027}

\section*{Literatur}\label{literatur}
\addcontentsline{toc}{section}{Literatur}

\protect\phantomsection\label{refs}
\begin{CSLReferences}{1}{0}
\vspace{1em}

\bibitem[\citeproctext]{ref-Badescu22Urban-Ge}
Bădescu, G. (2022). Urban Geopolitics in 'Ordinary' and 'Contested'
Cities: Perspectives from the European South-East. \emph{Geopolitics},
\emph{28}(5), 1757--1782.
\url{https://doi.org/10.1080/14650045.2022.2129010}

\bibitem[\citeproctext]{ref-Bieber20The-Rise}
Bieber, F. (2020). \emph{The Rise of Authoritarianism in the Western
Balkans}. London: Palgrave Macmillan.

\bibitem[\citeproctext]{ref-Bruff14The-Rise}
Bruff, I. (2014). The Rise of Authoritarian Neoliberalism.
\emph{Rethinking Marxism}, \emph{26}(1), 113--129.
\url{https://doi.org/10.1080/08935696.2013.843250}

\bibitem[\citeproctext]{ref-Camprag23Megaproj}
Čamprag, N. (2023). Megaprojekt ‚{Belgrade} {Waterfront}`: Neugestaltung
der Stadt durch autorit{ä}re Stadtentwicklung.
\emph{sub\textbackslash urban. zeitschrift f{ü}r kritische
stadtforschung}, \emph{11}(3/4), 43--72. Abgerufen von
\url{https://zeitschrift-suburban.de/sys/index.php/suburban/article/view/856}

\bibitem[\citeproctext]{ref-Camprag24Effects}
Čamprag, N. (2024). Effects and Consequences of Authoritarian Urbanism:
Large-Scale Waterfront Redevelopments in Belgrade, Zagreb, and Novi Sad.
\emph{Urban Planning}, \emph{9}. \url{https://doi.org/10.17645/up.7589}

\bibitem[\citeproctext]{ref-Cupac20The-anxi}
Cupać, J. (2020). The anxiety dilemma: Locating the Western Balkans in
the age of anxiety. \emph{Journal of Regional Security}, \emph{15}(1),
7--38. \url{https://doi.org/10.5937/jrs15-25359}

\bibitem[\citeproctext]{ref-Darieva26Connecti}
Darieva, T., Krüsmann, V., \& Langbein, J. (2026). \emph{Connectivity
and the Entrenchment of Authoritarianism in {Azerbaijan}, {Serbia} and
{Georgia}} (ZOiS Report No. 3/2026). Berlin: Zentrum f{ü}r Osteuropa-
und internationale Studien (ZOiS).

\bibitem[\citeproctext]{ref-Ejdus17Not-a-he}
Ejdus, F. (2017). {„Not a heap of stones``}: material environments and
ontological security in international relations. \emph{Cambridge Review
of International Affairs}, \emph{30}(1), 23--43.
\url{https://doi.org/10.1080/09557571.2016.1271310}

\bibitem[\citeproctext]{ref-Eror18How-Alek}
Eror, A. (2018, März 9). How {Aleksandar} {Vucic} Became {Europe}'s
Favorite Autocrat: The {EU} is undermining its credibility by choosing
stability over democracy in {Serbia}. Abgerufen 10. November 2026, von
\url{https://foreignpolicy.com/2018/03/09/how-aleksandar-vucic-became-europes-favorite-autocrat/}

\bibitem[\citeproctext]{ref-Grubbauer19Urban-me}
Grubbauer, M., \& Čamprag, N. (2019). Urban megaprojects, nation-state
politics and regulatory capitalism in {Central} and {Eastern} {Europe}:
The {Belgrade} {Waterfront} project. \emph{Urban Studies}, \emph{56}(4),
649--671. \url{https://doi.org/10.1177/0042098018757663}

\bibitem[\citeproctext]{ref-Piletic22The-role}
Piletić, A. (2022). The role of the urban scale in anchoring
authoritarian neoliberalism: a look at post-2012 neoliberalization in
{Belgrade}, {Serbia}. \emph{Globalizations}, \emph{19}(2), 285--300.
\url{https://doi.org/10.1080/14747731.2021.1882817}

\bibitem[\citeproctext]{ref-Zekovic23The-fina}
Zeković, S., Perić, A., \& Hadžić, M. (2023). The financialization of
{„the urban``} in the post-socialist {Serbia}: evidence from the
{Belgrade} {Waterfront} megaproject. \emph{Journal of Urban Affairs}.
\url{https://doi.org/10.1080/07352166.2023.2239962}

\end{CSLReferences}




\end{document}
